\paragraph{Lựa chọn đặc trưng}

Lựa chọn đặc trưng nhằm mục tiêu lựa chọn một tập con các đặc trưng quan trọng
nhất, giúp giảm độ phức tạp của mô hình, cải thiện khả năng tổng quát hóa và
giảm hiện tượng overfitting. Dưới đây là một số phương pháp phổ biến.

\subparagraph{Forward Stepwise Selection}
Phương pháp này bắt đầu với một mô hình rỗng (không có đặc trưng nào). Tại mỗi
bước, ta thêm vào đặc trưng giúp cải thiện hiệu năng mô hình nhiều nhất (ví
dụ: giảm sai số hoặc tăng độ chính xác).
\begin{itemize}
    \item Khởi đầu với tập đặc trưng rỗng.
    \item Tại mỗi bước, thử thêm từng đặc trưng chưa được chọn vào mô hình.
    \item Chọn đặc trưng làm cải thiện tốt nhất theo tiêu chí đánh giá.
    \item Lặp lại cho đến khi không còn cải thiện đáng kể hoặc đạt số lượng
    đặc trưng mong muốn.
\end{itemize}

\subparagraph{Backward Elimination}
Ngược lại với forward selection, phương pháp này bắt đầu với toàn bộ tập đặc
trưng và loại bỏ dần các đặc trưng kém quan trọng.
\begin{itemize}
    \item Khởi đầu với tất cả các đặc trưng.
    \item Tại mỗi bước, loại bỏ đặc trưng ít đóng góp nhất (ví dụ: có trọng số
    nhỏ hoặc không có ý nghĩa thống kê).
    \item Đánh giá lại mô hình sau khi loại bỏ.
    \item Lặp lại cho đến khi việc loại bỏ làm giảm hiệu năng đáng kể.
\end{itemize}

\subparagraph{Feature Selection bằng Lasso}
Phương pháp này dựa trên regularization với chuẩn $L_1$ (Lasso). Khi huấn
luyện mô hình với Lasso, một số hệ số $w_j$ sẽ bị đưa về đúng $0$.
\begin{itemize}
    \item Huấn luyện mô hình với regularization $L_1$.
    \item Các đặc trưng có hệ số $w_j = 0$ được xem là không quan trọng và bị
    loại bỏ.
    \item Các đặc trưng có hệ số khác $0$ được giữ lại.
\end{itemize}

Ưu điểm của phương pháp này là thực hiện feature selection một cách tự động
trong quá trình huấn luyện, đặc biệt hiệu quả khi số chiều lớn.
